\documentclass[a4paper]{article}
\usepackage[utf8]{inputenc}
\usepackage{hyperref}
\usepackage{amsmath}
\usepackage{amssymb}
\usepackage{amsfonts}
\usepackage[backend=bibtex]{biblatex}
\bibliography{biblio}

\title{Lattice Signature with Efficient Protocols}
\author{Bataille Alex, Clerbaux Coline\\Fallik Gabriel, Grasso Francesca, Tajani Mohamed}

\begin{document}
\maketitle
\section{Introduction}
Digital signatures are fundamental cryptographic primitives, allowing a user to authenticate a message and guarantee its integrity. For decades, 
most practical signature schemes have relied on classical hardness assumptions such as integer factorization or discrete logarithms. However, 
these assumptions are known to be broken by quantum algorithms, which motivates the search for alternative constructions secure in the post-quantum setting.
Lattice-based cryptography has emerged as one of the most promising frameworks for post-quantum security, as its underlying problems are believed to remain 
hard even for quantum computers. Beyond security, modern applications often require signatures to interact efficiently with zero-knowledge proof systems, 
for instance in privacy-preserving protocols. Designing lattice-based signatures that remain compatible with such proof systems is therefore crucial.
A first attempt in this direction was the construction of Libert, Ling, Mouhartem et al. (\cite{LLM+16}), which proposes a lattice-based signature scheme with 
efficient protocols. Although conceptually appealing, their construction suffers from extremely large proof sizes (several hundreds of megabytes) making it impractical.
The paper studied in this report (\cite{mainPaper}) revisits this problem and proposes a new lattice-based signature scheme specifically designed to interface smoothly 
with recent zero-knowledge argument systems, while reducing the size of the produced proofs. In the following, we analyse the main ideas behind 
this construction, highlight its improvements over previous work, and examine some of its technical components in more detail.

\section{Paper Contributions}
\subsection{Large Sizes Problem}
As mentioned in the introduction, the signature scheme from Libert et al. \cite{LLM+16} suffers from large sized proofs (30 to 40 times larger). In this section, we study why 
this previous scheme yields such enormous sizes. In \cite{LLM+16}, the signature scheme relies on a commitment $\mathbf{c} = \mathbf{D_0 r} + \mathbf{D_1 m}$. To 
prove knowledge of a valid signature on a hidden message, the prover must demonstrate knowledge of the opening $(\mathbf{r}, \mathbf{m})$ such that $\mathbf{c}$ is well formed.
More precisely, the verification algorithm operates on the binary decomposition of the commitment vector. In our case, $\mathbf{c} \in \mathbb{Z}_q^n$ so the 
protocol requires decomposing each coefficient $c_i$ into $l = \lceil \log_2 q\rceil$ bits. \\
This design choice induces 2 main inefficiencies : 
\begin{enumerate}
    \item Witness expansion : The witness size for the ZKAoK scales linearly with the number of bits to be proven. Instead of proving knowledge of $n$ elements
            in $Z_q^n$, the prover must prove knowledge of $n \cdot \log_2 q$ binary variables. 
    \item Constraint proliferation : To ensure the decomposition is valid, the ZK protocol must enforce the binary constraint for every bit. In frameworks 
            based on quadratic constraints, this requires adding $\mathcal{O}(n \log q)$ quadratic equations to the proof system. Furthermore, the security
            reduction in \cite{LLM+16} relies on embedding a trapdoor within the decomposition matrix. This introduces a norm growth factor $\propto \sqrt{n\log q}$
             in the extracted solution, forcing the scheme to adopt overly conservative parameters to maintain security margins. 
\end{enumerate}

\subsection{The Solution}
To mitigate this, we optimize the interactions between the commitment scheme implicitly used by privacy-enhancing protocols and the 
signature scheme itself. Instead of generating the public parameters of these 2 components independently from one another, the authors propose a unified 
generation, using a common matrix $\mathbf{A}$. This indeed reduces the size of the proofs since the size of the 
public key is smaller and it allows the commitment randomness $\mathbf{r}$ to be algebraically absorbed into the signature vector $\mathbf{v}$, eliminating
the need for a separate proof of opening. 
\section{Signature Scheme}
The signature scheme proposed by the paper is based on the one studied in \cite{LLM+16}, along with the improvements discussed in the previous section. 

\section{Technical Analysis}

\section{Conclusion}

\newpage
\printbibliography

\end{document}
